\documentclass[11pt]{article}

\usepackage[margin=1in]{geometry}
\usepackage{amsmath,amssymb,amsthm}
\usepackage{booktabs}
\usepackage{hyperref}
\usepackage{longtable}

\newtheorem{theorem}{Theorem}
\newtheorem{lemma}{Lemma}
\newtheorem{claim}{Claim}
\newtheorem{definition}{Definition}
\newtheorem{remark}{Remark}

\title{A Machine-Verified Certificate That No Conway 99-Graph\\Contains a Paley-Perfect Vertex}
\author{R230 review bundle (claim corrected after adversarial review)}
\date{July 1, 2026}

\begin{document}
\maketitle

\begin{abstract}
This note documents the R230 computational certificate for the following
statement: no strongly regular graph with parameters
\(\operatorname{srg}(99,14,1,2)\) contains a \emph{Paley-perfect} vertex ---
a vertex all 21 of whose rooted far-cell fibers induce a 4-cycle
(equivalently, a vertex at which the R204 forced-edge table holds; equivalently,
a vertex all of whose 21 rooted quadrilateral closures induce the Paley graph of
order 9).  The proof is computational: any such rooted configuration reduces
exactly to a 105-block \(S_4\) permutation system, which splits into 24
symmetry representatives, each of which is refuted by a CNF with an
independently checked DRAT proof of unsatisfiability.

This strengthens the closest published results: Makhnev (1988) and, independently,
Keramatipour (2023) proved that no \(\operatorname{srg}(99,14,1,2)\) satisfies
the corresponding \emph{global} hypothesis (every vertex Paley-perfect;
Makhnev's condition~(*), Keramatipour's ``Paley(9) pattern'').  The present
certificate refutes a single Paley-perfect vertex.  It does \emph{not} resolve
Conway's 99-graph problem: a hypothetical graph in which every vertex has at
least one non-\(C_4\) fiber remains unexcluded, and the conjectured direction
in the literature (Keramatipour, Conjecture~3.4.4) is that no Paley(9) subgraph
occurs at all.  An earlier version of this bundle overclaimed unconditional
nonexistence; the adversarial review that led to this correction is included in
the repository.
\end{abstract}

\section{Statement}

\begin{definition}[Rooted fibers]
Let \(\Gamma\) be a hypothetical \(\operatorname{srg}(99,14,1,2)\) and let
\(v\) be a vertex.  Since \(\lambda=1\), the neighborhood \(N(v)\) induces a
perfect matching \(7K_2\); write its 7 matched pairs as
\(\{p_0^i,p_1^i\}\), \(i=0,\dots,6\).  Every vertex at distance two from
\(v\) (a \emph{far} vertex) is adjacent to exactly \(\mu=2\) vertices of
\(N(v)\), never to both ends of one matched pair; hence far vertices are in
bijection with the \(\binom{14}{2}-7=84\) non-matched local pairs.  Group the
84 labels by the (unordered) pair of matched-pair indices they meet: this
yields \(\binom{7}{2}=21\) \emph{fibers} of four far vertices each.
\end{definition}

\begin{definition}[Paley-perfect vertex]\label{def:pp}
The vertex \(v\) is \emph{Paley-perfect} if every one of its 21 fibers induces
a 4-cycle in which two far vertices are adjacent exactly when their labels
share one local vertex.  Equivalently (Lemma~\ref{lem:equiv}), the full R204
forced-edge table holds at \(v\): fibers are such \(C_4\)'s \emph{and} far
vertices in fibers sharing exactly one matched-pair index are non-adjacent.
Equivalently, for each fiber the nine vertices consisting of \(v\), the four
local vertices met by the fiber, and the four fiber vertices induce
\(\operatorname{srg}(9,4,1,2)\), the Paley graph of order~9.
\end{definition}

\begin{theorem}[R230 certificate, corrected claim]\label{thm:main}
No \(\operatorname{srg}(99,14,1,2)\) contains a Paley-perfect vertex.
\end{theorem}

The proof is a checked computation whose artifacts (CNFs, DRAT proofs, audit
scripts, hashes, and logs) constitute the certificate; every artifact-level
claim below was independently re-verified during an adversarial review
(byte-identical CNF rebuilds; all 24 DRAT proofs re-checked with an
independently compiled \texttt{drat-trim}).

\begin{lemma}[Equivalences in Definition~\ref{def:pp}]\label{lem:equiv}
At a fixed root \(v\): all 21 fibers being \(C_4\)'s (adjacent exactly on
one shared label vertex) is equivalent to the full forced-edge table at
\(v\), and to all 21 nine-vertex fiber closures inducing Paley(9).
\end{lemma}

\begin{proof}[Proof sketch]
\(\lambda=1\) makes every neighborhood a perfect matching, so for each local
vertex \(\ell\) and far \(x\ni\ell\), the edge \(\{\ell,x\}\) lies in a unique
triangle whose third vertex is the unique neighbor of \(x\) whose label
contains \(\ell\).  If fibers are \(C_4\)'s, that slot is filled inside
\(x\)'s own fiber, so no neighbor of \(x\) in a fiber sharing one index can
contain \(\ell\); the \(\mu\)-count for the mate of \(\ell\) closes the
remaining (label-disjoint) intersecting case.  Conversely, given the
intersecting-fiber rule, the matching partner of \((\ell,w)\) within the
\(\ell\)-family must be \((\ell,\operatorname{mate}(w))\) --- the only
candidate not in an intersecting fiber --- which yields the \(C_4\); the
opposite-corner non-edge follows from \(\lambda=1\) since both shared
corners are already common neighbors.  The nine-vertex closure of a
\(C_4\)-fiber is 4-regular on 9 vertices with \(\lambda=1,\mu=2\), which is
uniquely Paley(9).
\end{proof}

\section{Reduction, given a Paley-perfect root}

Fix a hypothetical \(\Gamma\) with a Paley-perfect vertex \(v\) and root at
\(v\).  The forced-edge table now holds \emph{by hypothesis}.  The R204
reduction then shows the remaining free far-edge surface is exactly a
105-block \(S_4\) permutation system: each far vertex has far-degree 12, two
forced \(C_4\)-neighbors, and no neighbors in intersecting fibers, hence 10
free neighbors in the 10 disjoint fibers; the same-fiber
\(\lambda/\mu\)-equations force the four free-neighbor sets of a fiber to be
pairwise disjoint, and \(4\times 10 = 40 = 10\times 4\) forces every disjoint
\(4\times4\) fiber block to be a permutation matrix.  The compact R204
far-pair equations are exactly the rooted SRG equations: the clean-room
symbolic audit
\[
\texttt{source/root\_cell\_r204\_cleanroom\_symbolic\_audit.py}
\]
verifies the symbolic identity of the compact and full equations for all 3360
distinct-fiber far pairs (\texttt{ok=true},
\texttt{symbolicPairEquationsChecked=3360}, \texttt{symbolicMismatches=[]}).
Note what this audit does and does not do: it verifies the algebra
\emph{downstream} of the forced-edge table; the table itself enters
Theorem~\ref{thm:main} as the Paley-perfect hypothesis, not as a proved fact.

\section{Triangle representative split and coset cuts}

The R220 split fixes a matching triple of disjoint fibers and quotients the
\(24^3=13824\) triples of \(S_4\) blocks by the rooted symmetry action.  The
audit \texttt{root\_cell\_triangle\_orbit\_audit.py} verifies 24 orbits whose
sizes sum to exactly 13824 and that the abstract
\(D_8^3\rtimes S_3\) action coincides with genuine rooted relabelings.  The
adversarial review additionally re-ran the quotient on a second, disjoint
matching triple \(((0,2),(1,4),(3,6))\) and obtained the identical orbit
structure, and verified that the wreath group \(S_2\wr S_7\) acts transitively
on each far-pair class (evidence in
\texttt{scratchpad/second\_triangle/} and
\texttt{scratchpad/pair\_orbit\_check/}).

R229 adds, for each of the 105 intersecting fiber pairs, the exact
right-\(D_8\)-coset projection of the residual relative-permutation table (182
allowed rows of \(3^6=729\)); the audit
\texttt{root\_cell\_intersecting\_coset\_sat\_audit.py} checks the projection
against the source relation and the relative-permutation orientation over all
\(24\times24\) block pairs.  These cuts are exact consequences of the
hypothesis-plus-permutation structure.

\section{SAT certificate}

Each of the 24 representatives is encoded (seqcounter cardinality plus R229
cuts) and refuted: CaDiCaL 3.0.0 with \texttt{--no-binary} produced ASCII DRAT
proofs, each checked by \texttt{drat-trim}.  The summary
\texttt{artifacts/audit\_json/r229\_all24\_ascii\_drat\_checked\_summary.json}
records \texttt{reps=24}, \texttt{unsatCount=24}, \texttt{verifiedCount=24}
with SHA-256 hashes; CNF and DRAT bodies ship via Git LFS.  The adversarial
review independently rebuilt all 24 CNFs byte-identically from
\texttt{source/root\_cell\_permutation\_sat.py} (all manifest hashes match)
and re-verified all 24 DRAT proofs with an independently compiled
\texttt{drat-trim}.

\medskip
\noindent\textbf{Why Theorem~\ref{thm:main} follows.}  A graph with a
Paley-perfect vertex yields, after rooting there, a 105-block \(S_4\)
assignment satisfying all rooted equations; the R220 action moves its fixed
triangle triple into one of the 24 verified representatives; the R229 clauses
are implied, so the assignment would satisfy one of the 24 CNFs --- all of
which are UNSAT with independently verified proofs. \qed

\section{Relation to the literature}

The global version of the hypothesis is not new, and its refutation is not
new either:

\begin{itemize}
\item A.~A.~Makhnev, \emph{On strongly regular graphs with \(\lambda=1\)}
(Mat.\ Zametki \textbf{44}:5 (1988) 667--672; Math.\ Notes \textbf{44} (1988)
847--850), defines condition~(*) and proves (Theorem~2) that no
\(\operatorname{srg}(99,14,1,2)\) satisfying~(*) exists.  Under~(*), every
quadrilateral generates an induced \(\operatorname{srg}(9,4,1,2)\); (*) is
the ``every vertex Paley-perfect'' hypothesis.
\item A.~Keramatipour, \emph{Approaching the Conway-99 problem using SAT
solvers} (MPhil thesis, Cambridge, 2023; arXiv:2604.23037), Definition~12
(``Paley(9) pattern'') is precisely the fiber-\(C_4\) table at every vertex;
his Theorem~3.4.2 independently re-proves Makhnev's Theorem~2.  His
Conjecture~3.4.4 asserts that \emph{no} Paley(9) subgraph occurs in a
99-graph at all.
\item P.~G.~Cesarz and A.~J.~Woldar, \emph{On the automorphism group of a
putative Conway 99-graph} (Algebraic Combinatorics; arXiv:2308.02978), use
the same rooted labeling; their unconditional Lemma~4.8 is the coordinate
balance underlying the fiber analysis, and their Remark~5.6 records that even
the existence of a single Paley(9) inside a 99-graph is unresolved.
\item S.~Lou and M.~Murin (MIT PRIMES 2014) prove the closure lemma that
Paley(9)-minus-an-edge forces the missing edge; R.~Reimbayev
(arXiv:2508.03377) gives a six-vertex census whose free parameter \(n_3\)
vanishes exactly under the global hypothesis, citing Makhnev for the fatal
consequence.
\end{itemize}

Theorem~\ref{thm:main} strengthens Makhnev/Keramatipour from ``every vertex''
to ``some vertex'', and machine-certifies the result.  To our knowledge the
single-vertex statement is new; it is also an unconditional step toward
Keramatipour's Conjecture~3.4.4, which would imply it.

\section{What remains open}

Conway's problem itself.  A hypothetical
\(\operatorname{srg}(99,14,1,2)\) in which every vertex has at least one
non-\(C_4\) fiber is untouched by this certificate.  Deciding the
corresponding ``flip'' SAT instances in either direction would resolve the
problem outright (UNSAT for the intersecting-fiber flips would prove the
forced table, contradicting Theorem~\ref{thm:main}'s hypothesis class being
empty \emph{via} Makhnev's route; SAT would exhibit the graph), which is why
no such certificate should be expected cheaply.

\section{Replay}

\begin{enumerate}
\item \texttt{python scripts/clean\_room\_replay.py} --- R204 symbolic audit,
metadata audit, hash/log certificate audit.
\item Re-run the R204/R220/R229 audits from \texttt{source/}.
\item Rebuild the 24 CNFs (\texttt{scripts/rebuild\_all24\_cnf.ps1}) and
compare SHA-256 hashes against
\texttt{artifacts/large\_artifacts\_manifest.csv}.
\item Re-check every DRAT proof with an independently built
\texttt{drat-trim}.
\end{enumerate}

\section{Provenance of the correction}

An earlier revision of this document claimed unconditional nonexistence,
treating the forced-edge table as established.  An adversarial review
(2026-07-01, included in the repository) identified the table as an unproven
load-bearing assumption, reduced it to the single intersecting-fiber kernel
lemma, and a subsequent literature search located Makhnev's and
Keramatipour's refutations of the global hypothesis --- establishing both
that the table cannot be ``certified'' short of resolving Conway's problem,
and that the correct scope of the present certificate is
Theorem~\ref{thm:main}.
\end{document}
